\section{From gravitational antenna theory to a two-ended passive bound}
\label{sec:context}

Classical gravitational-antenna theory already contains much of the single-interface physics needed here. For compact resonant mechanical antennas, Hirakawa, Narihara, and Fujimoto developed a common eigenmode description of gravitational emission and reception, defined a mode coupling resource from the dynamic mass quadrupole, proved emission--reception reciprocity, and derived the corresponding directivity patterns \cite{HirakawaNariharaFujimoto1976}. Their short-pulse response is independent of mechanical quality factor and is controlled by the incident spectral energy density, making the oscillator-strength character of the integrated response explicit. Resonant-bar calculations likewise show that frequency-integrated gravitational absorption is controlled by oscillator strength rather than growing indefinitely with $Q$ \cite{PaikWagoner1976}. Material-response descriptions later express gravitational absorption through dynamical elastic or stress susceptibilities \cite{SrivastavaWidomPizzella2003}. We take compact antenna reciprocity, quadrupole-controlled oscillator strength, and $Q$-independent integrated response as prior physics, not as new results.

Complete laboratory generator--detector calculations are also historical. Grishchuk and Sazhin analyzed an actively driven electromagnetic gravitational-wave source together with a resonant electromagnetic receiver and obtained architecture-specific constraints on the complete system \cite{GrishchukSazhin1975}. Their extended, driven standing-wave architecture is outside the passive compact wave-zone class studied below, but it rules out any novelty claim based merely on retaining a source and receiver in one gravitational calculation.

A related efficiency--bandwidth viewpoint has become central in quantum transduction. A useful transducer is characterized not only by on-resonance conversion efficiency but by the frequency-dependent quantum channel, including bandwidth and added noise. Continuous-time quantum capacities are naturally written as frequency integrals of per-frequency channel capacities \cite{WangLiJiang2022}, while passive cavity transducers obey efficiency--bandwidth constraints that can be modified only by introducing additional physical resources such as nonclassical assistance \cite{ShiZhuang2024}. These developments motivate a narrower gravitational question:
\begin{quote}
For two separated compact passive matter interfaces connected by direct propagating gravity, can the frequency-integrated local-port-to-local-port coherent transfer be bounded by microscopic gravitational oscillator-strength resources at both endpoints and by the normalized free-space TT channel between them?
\end{quote}

This question differs from either a single gravitational antenna or an architecture-specific generator--detector feasibility calculation. We seek a many-mode passive cut set in which the source gravitational coupling resource, the propagating TT channel, and the receiver gravitational coupling resource are bounded separately and then closed microscopically. Recent work has also formulated gravity-mediated quantum channels, gravitational quantum communication benchmarks, and graviton-sensitive receivers in increasingly explicit ways \cite{MariEtAl2026,ToccaceloAndersenBrask2025,TobarEtAl2024}. Our purpose here is specifically the passive integrated resource inequality, not the existence of gravitational emission, reception, or source--receiver coupling itself.

The central quantity is the frequency-integrated coherent transmission. Let $T(\omega)$ be the matrix transfer function from selected local input channels of the source device to selected local output channels of the receiver device, after all unobserved passive ports have been retained in the physical model. We define
\begin{equation}
\boxed{
\Gamma_{\rm coh}
\equiv
\frac{1}{2\pi}
\int_{\mathcal B}\dd\omega\,
\Tr[T^\dagger(\omega)T(\omega)].}
\label{eq:Gammacoh}
\end{equation}
For a scalar channel this is simply the area under its transmissivity spectrum. It has units of inverse time and contains no arbitrary full-width-at-half-maximum convention. It is not itself a Shannon or quantum capacity; information-theoretic consequences will be stated only after the physical response bound is established.

We derive the result in three steps. First, established passive linear-system identities are used to bound the integrated useful transfer by the smaller gravitational coupling resource of the two endpoints. Second, for compact nonrelativistic linear matter, the mass-quadrupole energy-weighted sum rule bounds the cumulative resource independently of how it is distributed among passive resonances. Third, the transverse-traceless radiation projector bounds the largest normalized propagation channel between compact quadrupoles in the wave zone. The resulting narrowband theorem is
\begin{equation}
\boxed{
\Gamma_{\rm coh}
\lesssim
\frac{25G\omega^2}{12c^3R^2}
\min\!\left(\langle I_A\rangle,\langle I_B\rangle\right),}
\label{eq:headline}
\end{equation}
where $I_A$ and $I_B$ are the internal mass inertia moments about the endpoint centers of mass. The theorem is restricted to the compact, passive, nonrelativistic, linear-bosonic, direct wave-zone class stated explicitly below. In particular, it is not a bound on active transducers, extended phased apertures, near-field gravity, higher multipoles, or relativistic sources.
